\section{Finite volumes and conventions}\label{setup:conventions}

Throughout the local and multiscale arguments, fix
\begin{equation}\label{setup:fixed-energy}
 h>0,\qquad -4-h<E<4+h.
\end{equation}
All constants and auxiliary choices may depend on this pair. Their
dependence on the varying geometric scales, exterior data and masks is
specified in each assertion. Uniformity in $E$ is neither assumed nor
needed for \cref{main:theorem}.

For integer $l\ge0$, put $b_l=3^l$. The global finite volume is the square
torus
\[
 \mathbb T_N^2=(\Z/b_N\Z)^2,
\]
with the periodic nearest-neighbor adjacency and independent uniform
potentials on its sites. Its area is $b_N^2$. Distances are shortest
coordinate distances on this torus, using the maximum norm unless a
different norm is displayed. We write
\[
 \Lambda_r(x)=\{y:\abs{y-x}_\infty\le r\}.
\]
All local constructions are used with the torus sufficiently large to
contain the required neighborhoods without ambiguity. Dirichlet
restrictions mean deletion of the hopping bonds leaving the indicated
physical site set.

The physical disorder is always the product law of the true variables
$v_x$. Auxiliary selectors, acceptance uniforms, offsets and thinning
marks will be introduced explicitly. Conditioning on a rare event can
change the law within one cell; it never changes this unconditional
physical law. An expectation refers to the sampling and conditioning
specified in its statement.

Site and port spaces are complexified finite-dimensional Euclidean
spaces. A superscript $\top$ denotes transpose, not adjoint; a superscript
$*$ denotes adjoint. Bases used for coordinate projections and their
subspaces are real orthonormal bases. We use
\[
 \operatorname{Im}D=\frac{D-D^*}{2i},\qquad
 \operatorname{Re}D=\frac{D+D^*}{2}.
\]
Unless otherwise indicated, a norm is the Euclidean or induced operator
norm. Hilbert--Schmidt and trace norms carry subscripts $\mathrm{HS}$
and $1$. Coordinate projections may be identified with their ranges.
Symbols used for local blocks are reset when the local matrix is defined.

A real reference potential may replace a true potential only at a site
carrying an open cell port. Closing that port restores the true value.
The matrix $H^\#$ denotes the real symmetric operator with the current
reference values. Artificial probes, introduced in the next section, do
not alter real reference values. Several ports may share a physical site;
their columns are kept distinct.

In bounds for a cell of size $L$, the notation $\operatorname{poly}(L)$
means a quantity bounded by $C(1+L)^C$ for fixed constants. The constants
may depend on the previously fixed cell-design parameters, including
the corner-box width exponent, but not on the exterior matrix, exterior
volume, exterior source, or allowed probe mask. The ordering of design
parameters is part of each construction. Expressions
$\exp(L^{O(1)})$ have the corresponding interpretation. A statement
requiring sufficiently large $L$ is applied only after its fixed design
parameters have been chosen.

We will repeatedly distinguish the following orders of choice. The
disorder and energy are fixed first. Local design constants are fixed
before the cell scale is enlarged. For global limits a starting level
$n$ is fixed before observation and approximation constants are chosen;
a later level $m$ can then be increased to exceed all those thresholds,
and the torus level $N$ is increased last. No rate uniform in $n$ is
implicit in this convention.
